% ===================================================================
%  BẢNG Số 1 — Trường học. --- Trường trung học. --- Lớp học.
%
%  Ceci n'est PAS une transcription : c'est une traduction, faite en
%  2026, en vietnamien standard d'aujourd'hui. D'ou trois differences
%  avec texto/io et texto/fr, et trois seulement :
%
%   * pas de \nl ni de \cc — la traduction ne suit aucune ligne
%     imprimee, et les lignes de ce fichier ne sont que du confort ;
%   * pas de \begin{VUpage} — elle n'a ni feuillet ni folio, et la
%     page de lecture ne lui en affiche pas ;
%   * le gras se pose sur le TERME seul, ponctuation dehors.
%
%  Tout le reste est identique, et doit l'etre : les cles « %%K » sont
%  celles de texto/io, macro pour macro pour l'apparat, et les renvois
%  (n) visent les MEMES objets de la planche.
%
%  LES DIX-SEPT LANGUES DE LA COMMUNAUTE IDISTE SONT CLOSES ; ON
%  REPREND ICI LE PROGRAMME DES VINGT-NEUF LANGUES D'ETHNOLOGUE, la ou
%  il s'etait arrete apres le turc. Le vietnamien vient le premier des
%  douze qui restent, avec quatre-vingt-six millions de locuteurs
%  premiers.
%
%  UNE DECISION DE COMPOSITION, ECRITE ICI POUR QU'ON NE LA CHERCHE
%  PAS AILLEURS. Le vietnamien ecrit « BẢNG Số 1 ». Le mot ne se
%  confond avec aucun autre membre du motif NUMERO_TAB — nulle autre
%  colonne n'ecrit un A a crochet — mais il a chez lui deux autres
%  emplois, et ils sont sur cette page meme : « bảng đen » est LE
%  TABLEAU NOIR de l'alinea 2, « bảng treo tường » LES TABLEAUX
%  MURAUX de l'alinea 2 aussi, et « bảng con » L'ARDOISE du petit
%  Emilô a l'alinea 12. La colonne s'appelle donc du nom de l'objet
%  qu'elle decrit, quatre fois dans le meme tableau. Rien de tout
%  cela ne passe le motif : la comparaison est sensible a la casse,
%  et seul l'apparat compose en capitales. Le cablage a recu
%  « BẢNG », « loạt » et « cảnh ».
%
%  ET LE VIETNAMIEN RANGE SON ORDINAL APRES LE NOM — « LOẠT THỨ
%  NHẤT », la serie ordre premier ; « Cảnh thứ nhất », scene ordre
%  premier. C'est la quatrieme langue de la serie a le faire, apres
%  le hindi, le bengali et les deux pendjabis, et comme eux il lui
%  faut son propre membre dans SERIO et dans CENO : le groupe commun
%  attend l'ordinal devant.
%
%  CE QUE LA PAGE MONTRE : TROIS COUCHES, ET CHACUNE NOMME UNE SEULE
%  ESPECE DE CHOSE.
%
%   1. LES DISCIPLINES SONT TOUTES SINO-VIETNAMIENNES. « hình học »
%      la geometrie, « số học » l'arithmetique, « vật lý » la
%      physique, « hóa học » la chimie, « động vật học » la zoologie,
%      « thực vật học » la botanique, « địa chất học » la geologie,
%      « lịch sử » l'histoire : pas une qui ne vienne du chinois. Ce
%      sont les mots de l'examen mandarinal, reemployes tels quels
%      pour le programme francais.
%
%   2. LES OBJETS QU'ON TIENT EN MAIN SONT VIETNAMIENS. « vở » le
%      cahier, « giấy » le papier, « mực » l'encre, « thước » la
%      regle, « phấn » la craie, « bảng » le tableau, « ghế » la
%      chaise, « bàn » la table, « tủ » l'armoire. Aucun n'a eu
%      besoin d'etre emprunte : ils etaient la avant l'ecole.
%
%   3. ET CE QUI EST ARRIVE DANS UNE CAISSE EST FRANCAIS. « com-pa »
%      le compas, « pi-a-nô » le piano, « gam » la gamme, et les sept
%      syllabes du solfege — đô, rê, mi, fa, son, la, si.
%
%  LES DEUX TABLEAUX MURAUX DE L'ALINEA 19 LE DISENT EN UN COUP
%  D'OEIL, cote a cote sur la meme cloison. Celui des sciences est
%  entierement chinois — vật lý, hóa học, động vật học, thực vật
%  học, địa chất học, lịch sử. Celui du systeme metrique est
%  entierement francais — « mét », « lít », « đề-ca-lít »,
%  « đề-xi-lít », « mét khối ». Deux empires sur le meme mur, et
%  l'ecolier lit les deux sans savoir qu'il change de langue. Le
%  tableau 1 estonien avait trouve le mobilier a l'ouest et le livre
%  a l'est ; ici la ligne ne passe pas entre deux marchands mais
%  entre L'ABSTRAIT ET LE MANIABLE.
%
%  ET LA NOTE (*) RECOIT UNE HUITIEME REPONSE. Guignon enumere
%  « Giovanni, Jean, Johann, Jan, Hans, John, Ivan » et demande qu'on
%  s'en tienne au nominatif latin. Le vietnamien en ajoute un
%  huitieme, « Gioan », et il n'est venu ni du latin ni du francais :
%  il est venu du PORTUGAIS, avec les jesuites de la mission, un
%  siecle et demi avant la colonie. Toute la colonne porte cette
%  marque, car le quốc ngữ lui-meme est une orthographe portugaise —
%  le « nh », le « gi », le « ch » final. Les noms de la page sont
%  donc ceux de la mission : Gioan, Phêrô, Phaolô, Giacôbê,
%  Stêphanô, Carôlô, Henricô, Antôn, Phanxicô. Les noms latins de la
%  note ne bougent pas : c'est d'eux qu'elle parle.
%
%  ENFIN, LE VIETNAMIEN N'A PAS DE GROUPE DE CONSONNES A L'INITIALE,
%  et il ne garde donc aucun mot etranger d'une seule piece : il le
%  coupe en syllabes et les relie d'un trait — « com-pa »,
%  « đề-ca-lít », « pi-a-nô », « Cơristianô ». C'est pourquoi cette
%  colonne porte des traits d'union la ou les seize autres n'en ont
%  pas : ils ne joignent pas des mots, ils separent des syllabes.
%
%  « postigo » (t01, renvoi 78) EST UN VASISTAS, ET LA PLANCHE LE DIT.
%  Le mot a ete verifie sur la gravure en 2026, non sur les
%  dictionnaires : la verriere de la classe a deux rangees de carreaux,
%  et dans la rangee HAUTE deux petits vantaux sont graves battants,
%  gonds sur un montant vertical, ouverts vers l'interieur. C'est le
%  seul ouvrant de la verriere, et il ne sert qu'a aerer.
%
%  D'OU CE QU'IL N'EST PAS, et ou seize colonnes s'etaient trompees :
%  ni une lucarne de toit (it. « lucernari »), ni un volet (cs, lt, lb),
%  ni un meneau ni un montant de chassis (ar, ca, es, oc), ni une
%  imposte fixe, ni le battant de la grande fenetre (bn), ni le
%  « ranma » de la cloison japonaise. Et « framuga » russe nomme
%  l'imposte entiere, quand le francais nomme le VANTAIL : c'est
%  « fortochka » qu'il fallait, comme l'ukrainien avait deja
%  « kvatyrka ».
% ===================================================================

%%K t01-apar-0 apar
\VUsaut{2.0mm}
\VUcentre{12.6pt}{120}{SÁCH GIẢI THÍCH}
\VUsaut{3.2mm}
\VUcentre{8.4pt}{160}{VỀ BỘ}
\VUsaut{2.6mm}
\VUtitre{15.6pt}{0}{Bảng trợ giúp Delmas}
\VUsaut{3.4mm}
\VUornamento{9pt}
\VUsaut{5.0mm}
\VUcentre{13.2pt}{90}{LOẠT THỨ NHẤT}
\VUsaut{3.0mm}
\VUfilet{16mm}
\VUsaut{5.6mm}

%%K t01-apar-1 apar
\VUtitre{15.0pt}{60}{BẢNG Số 1}
\VUsaut{1.4mm}
\VUtitre{11.4pt}{0}{Trường học. --- Trường trung học. --- Lớp học.}
\VUsaut{2.2mm}
\VUfilet{20mm}
\VUsaut{6.4mm}

%%K t01-tit-1 sub
\VUpk{10.2pt}{40}{Mô tả tổng quát.}
\VUsaut{3.0mm}

%%K t01-01-1 p
1. --- Bảng này trình bày \VUgras{lớp học} của một \VUgras{trường trung học}. Trong lớp
học này có tám cái \VUgras{bàn} \textsuperscript{(18)} và tám cái \VUgras{ghế dài}
\textsuperscript{(32)}. Trên mỗi ghế dài có ba học sinh ngồi. \VUgras{Thầy giáo}
\textsuperscript{(7)} ngồi trên \VUgras{ghế} \textsuperscript{(8)} sau \VUgras{bục giảng}
\textsuperscript{(6)} của mình.

%%K t01-02-1 p
2. --- Bên trái bục giảng có một \VUgras{tủ} cửa kính \textsuperscript{(15)}. Bên phải
là \VUgras{bảng đen} \textsuperscript{(1)} với \VUgras{giẻ lau} \textsuperscript{(2)} và
\VUgras{phấn} \textsuperscript{(3)}.

%%K t01-02-2 p
Phía trên bục giảng treo các \VUgras{bảng treo tường} \textsuperscript{(9, 11, 12)} và
\VUgras{bản đồ} \textsuperscript{(10)}.

%%K t01-03-1 p
3. --- Không phải học sinh nào cũng ngồi đúng \VUgras{chỗ} \textsuperscript{(17)} của
mình. Gioan \textsuperscript{(4)} đứng bên bảng; em cầm giẻ lau và xóa các
\VUgras{hình hình học}.

%%K t01-03-2 p
Phêrô đứng gần bục giảng; em thu các \VUgras{bài viết} \textsuperscript{(5)} và mang lên
cho thầy giáo.

%%K t01-04-1 p
4. --- Thầy giáo này là thầy dạy \VUgras{sinh ngữ}. Thầy dạy học sinh nói, đọc và viết
các thứ tiếng nước ngoài (\VUgras{tiếng Đức, tiếng Anh, tiếng Tây Ban Nha, tiếng Ý,
tiếng Nga} hoặc \VUgras{tiếng Pháp}).

%%K t01-05-1 p
5. --- Lớp học rất rộng rãi và rất sáng sủa. Ở bức tường cuối có một \VUgras{cửa sổ}
rộng, hai \VUgras{cửa sổ thông gió} \textsuperscript{(78)} và một \VUgras{cửa kính} để
thông khí và lấy ánh sáng.

%%K t01-05-2 p
Lớp học này không lạnh. Ở giữa có một \VUgras{lò sưởi} \textsuperscript{(46)} rất tốt
với \VUgras{ống khói} \textsuperscript{(45)} để sưởi ấm.

%%K t01-05-3 p
Trên tường có gắn những \VUgras{mắc áo} \textsuperscript{(58)}; trên đó treo mũ và áo
khoác của học sinh.

%%K t01-tit-2 sub
\VUsaut{5.2mm}
\VUpk{10.2pt}{40}{Sân chơi.}
\VUsaut{3.6mm}

%%K t01-06-1 p
6. --- \VUgras{Đồng hồ} \textsuperscript{(63)} chỉ chín giờ năm phút.

%%K t01-06-2 p
\VUgras{Giờ ra chơi} \textsuperscript{(76)} vừa kết thúc và giờ học bắt đầu lại. Giờ ra
chơi diễn ra ngoài \VUgras{sân} \textsuperscript{(75)}. Ta thấy vài học sinh đang chơi.
\VUgras{Giám thị} \textsuperscript{(74)} trông coi các em.

%%K t01-07-1 p
7. --- Ngoài sân ta còn thấy những người khác: \VUgras{thanh tra} \textsuperscript{(73)},
\VUgras{hiệu trưởng} \textsuperscript{(71)} và \VUgras{giám học} \textsuperscript{(72)}.

%%K t01-07-2 p
Thanh tra kiểm tra các trường học, hiệu trưởng điều hành trường trung học, giám học
trông coi việc học.

%%K t01-07-3 p
Người thứ tư là \VUgras{người gác cổng} \textsuperscript{(77)}. Ông cắp dưới nách quyển
sổ để ghi những người vắng mặt.

%%K t01-08-1 p
8. --- Ở cuối sân có \VUgras{dụng cụ thể dục} \textsuperscript{(70)} và
\VUgras{xưởng thủ công} \textsuperscript{(69)}.

%%K t01-08-2 p
Ở bên phải bảng ta bắt đầu thấy các \VUgras{phòng nhạc} và \VUgras{phòng vẽ}.

%%K t01-noto-1 noto
\VUnotes{143.2mm}{%
(*) Đối với \textit{tên thánh}, nên viết thêm cả dạng La-tinh. Đó là cách duy nhất để
tránh vô số hình thức riêng. Nếu không, biết chọn thế nào giữa \textit{Giovanni, Jean,}
\textit{Johann, Jan, Hans, John, Ivan} v.v.? Chẳng lẽ ta lại làm riêng một cuốn từ điển
tên thánh? Hãy lấy ở tiếng La-tinh dạng chủ cách của chúng mà nói: \VUgras{Ioannes,
Ioanna; Iulius, Iulia; Iakobus; Andreas; Lukas.} (Ngữ pháp đầy đủ).%
}

%%K t01-tit-3 sub
\VUpk{10.2pt}{40}{Mô tả chi tiết.}
\VUsaut{2.4mm}

%%K t01-tit-4 sub
\VUpk{10.2pt}{40}{Học sinh giỏi.}
\VUsaut{3.0mm}

%%K t01-09-1 p
9. --- Ở dãy ghế đầu bên phải bục giảng ngồi \VUgras{Henricô} (*), \VUgras{Stêphanô} và
\VUgras{Carôlô}.

%%K t01-09-2 p
Henricô rất chăm chú và siêng năng; em nghe thầy giảng. Trên bàn em có \VUgras{vở}
\textsuperscript{(28)}, \VUgras{sách} \textsuperscript{(29)}, \VUgras{từ điển}
\textsuperscript{(30)} và \VUgras{hộp bút} \textsuperscript{(31)}. Sách của em có nhiều
\VUgras{trang}; mỗi chương có mấy \VUgras{đoạn} và mỗi \VUgras{dòng} có nhiều
\VUgras{từ}.

%%K t01-09-4 p
Bạn ngồi cạnh em, Stêphanô \textsuperscript{(24)}, rất cần mẫn. Em ghi vào vở bài học
cho buổi sau. Em có \VUgras{nét chữ} đẹp. Sách của em đóng lại. Ta chỉ thấy
\VUgras{bìa} \textsuperscript{(27)}. Bên cạnh em có \VUgras{com-pa}
\textsuperscript{(26)}.

%%K t01-09-5 p
Carôlô \textsuperscript{(23)} cầm sách trên tay; em mở ra để ôn lại bài. Như mọi học
sinh, trước mặt em có \VUgras{lọ mực} \textsuperscript{(25)}. Em rất vâng lời.

%%K t01-10-1 p
10. --- Ở bàn bên cạnh ngồi \VUgras{Ludovicô, Antôn} và \VUgras{Phaolô}. Ludovicô đang
trả \VUgras{bài học} \textsuperscript{(22)} và trả lời các \VUgras{câu hỏi} của thầy.
Antôn \textsuperscript{(21)} rất thiếu chú ý; đôi khi thầy còn phạt em. Phaolô
\textsuperscript{(20)} là một \VUgras{người bạn} tốt, hiền lành và sẵn lòng giúp đỡ. Em
rất chăm chú.

%%K t01-tit-5 sub
\VUsaut{4.2mm}
\VUpk{10.2pt}{40}{Học sinh kém.}
\VUsaut{3.2mm}

%%K t01-11-1 p
11. --- Sau lưng Henricô ngồi \VUgras{Giorgiô} \textsuperscript{(38)}. Em không phải là
đứa trẻ hư, nhưng em \VUgras{hay nói chuyện}, thiếu chú ý và không chịu ngồi yên.
\VUgras{Tính nhẹ dạ} và \VUgras{sự không vâng lời} của em gây \VUgras{lộn xộn} trong
lớp, và em thường phải nhận \VUgras{lời cảnh cáo} nghiêm khắc. \VUgras{Vở nháp}
\textsuperscript{(35)} của em bẩn thỉu, em làm \VUgras{vết mực} lên đó bằng
\VUgras{ngòi bút} \textsuperscript{(34)}, và vì thế em luôn phải dùng \VUgras{cục tẩy}
\textsuperscript{(36)}; \VUgras{cặp sách} \textsuperscript{(33)} của em rách vì em rất
cẩu thả. Tại sao em lại mang \VUgras{ống đựng bút} \textsuperscript{(37)} vào giờ sinh
ngữ?

%%K t01-11-2 p
Bạn bên cạnh em, Étmunđô \textsuperscript{(39)}, không ưa em. Em này quả có hơi lười,
nhưng em im lặng và ngồi yên. Em không nói chuyện; chỉ có điều thay vì nghe giảng, em
thích đọc những \VUgras{truyện} trong \VUgras{sách tập đọc} của mình.

%%K t01-11-3 p
Học sinh thứ ba của dãy ghế này là \VUgras{Ludovicô} \textsuperscript{(40)}. Em cần mẫn.
Em viết trên \VUgras{tờ giấy} trắng. Trước mặt mình, em đã đặt \VUgras{giấy thấm}
\textsuperscript{(41)}.

%%K t01-12-1 p
12. --- Học sinh của dãy ghế thứ hai, bên trái thầy giáo, còn rất nhỏ. Em đầu tiên, bé
\VUgras{Emilô}, chưa biết viết cho thạo; em mang theo \VUgras{bảng con}
\textsuperscript{(42)}, \VUgras{bút đá} và \VUgras{bọt biển} \textsuperscript{(43)} của
mình.

%%K t01-12-2 p
Ngồi cạnh em là \VUgras{Augustô} và \VUgras{Bênađô} \textsuperscript{(44)}. Em sau này
học tốt, nhưng \VUgras{quần áo} của em rất luộm thuộm.

%%K t01-13-1 p
13. --- Sau lưng các em ngồi ba \VUgras{đứa lười. Albêtô} không học bài.
\VUgras{Giacôbê} \textsuperscript{(47)} ngủ thay vì nghe giảng. \VUgras{Sự lười biếng}
của em mang lại cho em \VUgras{lời trách mắng} và cả \VUgras{hình phạt}. Sau cùng,
\VUgras{Cơristianô} \textsuperscript{(48)} quá nhẹ dạ. Em \VUgras{cư xử} không tốt và
chẳng cố gắng gì để sửa mình.

%%K t01-14-1 p
14. --- Trái lại, những bạn ngồi cạnh em cư xử tốt. \VUgras{Ernestô}
\textsuperscript{(51)} yêu \VUgras{trật tự} và \VUgras{sự ngăn nắp}; em luôn dùng
\VUgras{thước kẻ} \textsuperscript{(50)} và \VUgras{bút chì} \textsuperscript{(49)} của
mình. Em là \VUgras{học sinh ngoại trú}.

%%K t01-14-2 p
\VUgras{Phanxicô} \textsuperscript{(52)} là \VUgras{học sinh nội trú}; em mặc
\VUgras{đồng phục} của trường; em ăn và ngủ ở đó. Em là một \VUgras{người bạn} tốt.

%%K t01-14-3 p
Môrixô gọt \VUgras{bút chì} bằng \VUgras{dao nhíp} \textsuperscript{(56)} của mình; em
rất cẩn thận.

%%K t01-14-4 p
Em mang theo tất cả sách vở của mình, cuốn \VUgras{ngữ pháp} \textsuperscript{(55)},
\VUgras{sổ tay} \textsuperscript{(54)} và cả \VUgras{tấm lót viết}
\textsuperscript{(53)}. \VUgras{Túi sách} \textsuperscript{(57)} của em nằm dưới sàn,
phía sau em.

%%K t01-14-5 p
Hai cái bàn cuối lớp là chỗ của những \VUgras{học sinh giỏi} \textsuperscript{(19)}.

%%K t01-tit-6 sub
\VUsaut{4.6mm}
\VUpk{10.2pt}{40}{Vẽ và âm nhạc.}
\VUsaut{3.4mm}

%%K t01-15-1 p
15. --- Trong \VUgras{lớp vẽ} có những \VUgras{mẫu vẽ} đẹp gắn trên tường và
\VUgras{đèn} \textsuperscript{(62)} có \VUgras{chao đèn}, treo trên trần.
\VUgras{Thầy giáo} \textsuperscript{(60)} rất kiên nhẫn; thầy sửa các \VUgras{bức vẽ}
\textsuperscript{(61)} của học sinh và dạy các em vẽ \VUgras{tượng bán thân}
\textsuperscript{(59)} theo mẫu thật.

%%K t01-16-1 p
16. --- \VUgras{Lớp nhạc} có vẻ rất thú vị. Ở đó người ta học đọc \VUgras{nhạc}, xướng
âm các \VUgras{nốt của gam} \textsuperscript{(67)} viết trên \VUgras{khuông nhạc} (đô,
rê, mi, fa, son, la, si\,=\,C. D. E. F. G. A. B.), học biết các \VUgras{khóa nhạc} và
hát những \VUgras{giai điệu} đẹp. Bản nhạc được đặt lên \VUgras{giá nhạc}
\textsuperscript{(65)}. Một học sinh \VUgras{chơi đàn dương cầm} \textsuperscript{(64)}
và \VUgras{thầy dạy nhạc} \textsuperscript{(66)} \VUgras{đánh nhịp}
\textsuperscript{(68)}.

%%K t01-17-1 p
17. --- Ta bắt đầu thấy một góc của \VUgras{xưởng thủ công} \textsuperscript{(69)}, nơi
các em trai có thể học bào gỗ và giũa sắt. Không phải trường trung học nào cũng có.

%%K t01-tit-7 sub
\VUsaut{5.0mm}
\VUpk{10.2pt}{40}{Lớp khoa học tự nhiên.}
\VUsaut{3.6mm}

%%K t01-18-1 p
18. --- Thầy dạy toán dạy học sinh \VUgras{hình học, số học, phép tính} và
\VUgras{hệ mét}. Trên bảng ta thấy những hình hình học chính: \VUgras{hình tròn}
\textsuperscript{(a)}, \VUgras{hình vuông} \textsuperscript{(b)}, \VUgras{hình tam giác}
\textsuperscript{(c)}, \VUgras{đường thẳng} \textsuperscript{(d)}.

%%K t01-18-2 p
Ta cũng thấy một \VUgras{phép tính}; đó không phải \VUgras{phép cộng}, cũng không phải
\VUgras{phép trừ} hay \VUgras{phép nhân}: đó là \VUgras{phép chia}
\textsuperscript{(e)}.

%%K t01-19-1 p
19. --- Trên bảng hệ mét ta phân biệt: \VUgras{mét} \textsuperscript{(a)},
\VUgras{cân} \textsuperscript{(b)} và các \VUgras{quả cân} của nó, \VUgras{lít}
\textsuperscript{(e)}, \VUgras{đề-ca-lít} \textsuperscript{(f)}, \VUgras{đề-xi-lít} và
\VUgras{mét khối} \textsuperscript{(d)}.

%%K t01-19-2 p
Học sinh cũng học \VUgras{khoa học tự nhiên} \textsuperscript{(11)} --- động vật học
\textsuperscript{(a)}, thực vật học \textsuperscript{(b)}, địa chất học
\textsuperscript{(c)} --- và \VUgras{lịch sử} \textsuperscript{(12)}.

%%K t01-19-3 p
Trong tủ có dụng cụ \VUgras{vật lý} \textsuperscript{(15)} và \VUgras{hóa học}
\textsuperscript{(16)}.

%%K t01-19-4 p
Trên nóc tủ có: \VUgras{quả cầu} \textsuperscript{(14)}, \VUgras{hình nón} và
\VUgras{quả địa cầu} \textsuperscript{(13)}.

%%K t01-19-5 p
Phía trên các bảng treo một tấm \VUgras{bản đồ} vẽ năm phần của thế giới:
\VUgras{châu Âu} \textsuperscript{(g)}, \VUgras{châu Á} \textsuperscript{(e)},
\VUgras{châu Phi} \textsuperscript{(f)}, \VUgras{châu Mỹ} \textsuperscript{(ab)} và
\VUgras{châu Đại Dương} \textsuperscript{(h)}, cùng hai đại dương lớn,
\VUgras{Thái Bình Dương} \textsuperscript{(c)} và \VUgras{Đại Tây Dương}
\textsuperscript{(d)}.
\VUsaut{6.0mm}
\VUfilet{22mm}
